\section{Latar Belakang}
Setelah melewati tahap lexical analysis, syntax analysis, dan semantic analysis, compiler telah memiliki representasi program yang valid secara sintaksis maupun semantik dalam bentuk decorated AST dan symbol table. Namun, decorated AST masih berbentuk pohon yang tidak dapat dieksekusi secara langsung oleh mesin. CPU dan memori komputer dirancang untuk mengeksekusi instruksi secara sekuensial, bukan dengan melompat-lompat dalam struktur pohon yang kompleks. Oleh karena itu, diperlukan sebuah jembatan yang meratakan decorated AST menjadi daftar instruksi sederhana dan sekuensial yang disebut intermediate code.

Pada Milestone 1, Arion Compiler telah memiliki lexer yang mengubah source code menjadi token stream. Pada Milestone 2, token tersebut diproses oleh Recursive Descent Parser untuk menghasilkan parse tree. Pada Milestone 3, parse tree dikonversi menjadi AST, dianalisis secara semantik, dan diperkaya dengan anotasi tipe data, scope, serta referensi ke symbol table. Hasil akhirnya adalah decorated AST beserta tiga tabel simbol \texttt{tab}, \texttt{btab}, dan \texttt{atab}.

Milestone 4 berfokus pada dua komponen terakhir dalam pipeline compiler Arion, yaitu intermediate code generation dan interpreter. Intermediate code generator bertugas mengubah decorated AST menjadi daftar instruksi standar yang mirip dengan bahasa mesin namun masih mudah dibaca manusia. Instruksi-instruksi ini kemudian dieksekusi oleh interpreter yang berperan sebagai virtual machine. Interpreter membaca instruksi satu per satu, mengelola memori melalui stack, mengevaluasi operasi matematika, mengatur alur kontrol program, dan menghasilkan output nyata ke layar.

\section{Rumusan Masalah}
Permasalahan utama pada Milestone 4 adalah bagaimana membangun intermediate code generator dan interpreter yang dapat memanfaatkan decorated AST hasil Milestone 3 untuk menghasilkan output final dari sebuah program Arion. Permasalahan tersebut diuraikan menjadi beberapa pertanyaan berikut.

\begin{enumerate}
    \item Bagaimana mengonversi decorated AST menjadi intermediate code berupa daftar instruksi sekuensial menggunakan pendekatan Three-Address Code?
    \item Bagaimana merancang instruction set yang mencakup operasi load, store, aritmatika, perbandingan, jump, call, dan return?
    \item Bagaimana mengimplementasikan stack machine untuk mengelola memori eksekusi, termasuk stack frame, static link, dynamic link, dan return address?
    \item Bagaimana melakukan translasi control flow (\texttt{if}, \texttt{while}, \texttt{for}, \texttt{repeat}, \texttt{case}) ke dalam instruksi \texttt{JMP} dan \texttt{JPC} dengan label?
    \item Bagaimana menangani kerentanan runtime seperti stack overflow, out-of-bounds access, invalid jump target, dan numerical overflow?
    \item Bagaimana menampilkan output intermediate code dan hasil eksekusi program secara terformat?
\end{enumerate}

\section{Tujuan}
Tujuan pengerjaan Milestone 4 adalah membangun tahap intermediate code generation dan interpreter pada Arion Compiler. Secara khusus, tujuan yang ingin dicapai adalah sebagai berikut.

\begin{enumerate}
    \item Membuat intermediate code generator yang dapat mengubah decorated AST menjadi daftar instruksi \texttt{LIT}, \texttt{LOD}, \texttt{STO}, \texttt{INT}, \texttt{JMP}, \texttt{JPC}, \texttt{OPR}, \texttt{CAL}, dan \texttt{RET}.
    \item Mengimplementasikan stack machine dengan siklus fetch-decode-execute untuk menjalankan intermediate code.
    \item Mengelola memori eksekusi melalui stack, termasuk alokasi frame, push/pop operand, dan pengaturan stack frame untuk procedure call.
    \item Menangani translasi control flow dari struktur pohon menjadi instruksi sekuensial dengan label dan jump.
    \item Menambahkan deteksi dan penanganan runtime vulnerabilities pada interpreter.
    \item Menghasilkan output berupa intermediate code dan output eksekusi program sebagai hasil akhir compiler.
\end{enumerate}

\section{Batasan Masalah}
Implementasi Milestone 4 dibatasi pada backend compiler Arion. Program tidak lagi melakukan analisis sintaksis atau semantik, karena validitas program diasumsikan sudah diperiksa oleh tahap sebelumnya. Intermediate code generator bekerja pada decorated AST yang dibangun dari parse tree hasil parser, sedangkan interpreter mengeksekusi intermediate code tanpa mempedulikan sintaks bahasa sumber.

Pengecekan runtime yang diimplementasikan berfokus pada stack overflow, stack underflow, stack smashing, out-of-bounds array access, invalid jump target, dan numerical overflow/underflow. Fitur yang belum didukung pada tahap ini meliputi function call dalam ekspresi (\texttt{FuncCall}), record field access dalam intermediate code, serta boolean operator \texttt{and}, \texttt{or}, dan \texttt{not} pada level code generation.
